\documentclass[11pt,a4paper]{article}

\usepackage[a4paper,margin=25mm]{geometry}
\usepackage{amsmath,amssymb}
\usepackage{array}
\usepackage{booktabs}
\usepackage{longtable}
\usepackage{enumitem}
\usepackage{xcolor}
\usepackage{hyperref}
\usepackage{fancyvrb}
\usepackage{microtype}

\hypersetup{
  colorlinks=true,
  linkcolor=blue!55!black,
  urlcolor=blue!55!black,
  citecolor=blue!55!black
}

\newcommand{\code}[1]{\ifmmode\text{\texttt{\detokenize{#1}}}\else\nolinkurl{#1}\fi}
\newcommand{\R}{\mathbb{R}}
\newcommand{\dd}{\,\mathrm{d}}
\newcommand{\norm}[1]{\left\lVert #1 \right\rVert}
\newcommand{\abs}[1]{\left\lvert #1 \right\rvert}
\newcommand{\oas}{\textsc{OAS}}

\title{Detailed Description of the \oas{} Nonlinear Solver}
\author{Code inspection note}
\date{May 8, 2026}

\begin{document}
\maketitle

\begin{abstract}
This document describes the nonlinear solution algorithm implemented in the
\oas{} solver, with emphasis on the steady-state nonlinear solver path used by
\code{SteadyStateNonLinearSolver}.  It follows the code from the model-level
time/load-step loop through the Newton-like nonlinear iterations, force
assembly, constraint reduction, linear solves, convergence checks, adaptive time
step control, restarts, and final state commit.  The goal is to make the
implementation readable from the solver variables themselves, so most formulas
are written next to the corresponding C++ member names.
\end{abstract}

\tableofcontents
\newpage

\section{Scope and Source Map}

The nonlinear implementation is spread across the model, solver, node, element,
and linear-algebra layers.  Table~\ref{tab:source-map} lists the files that are
most relevant for understanding the control flow.

\begin{longtable}{p{0.33\linewidth} p{0.60\linewidth}}
\caption{Main source files for the nonlinear solver.}
\label{tab:source-map}\\
\toprule
File & Role \\
\midrule
\endfirsthead
\toprule
File & Role \\
\midrule
\endhead
\path{src/solver/src/model.cpp} &
Owns the global analysis loop in \code{Model::solve()}.  It calls
\code{solver->solveStep()}, exports each accepted step, and reports step wall
time. \\

\path{src/solver/src/solver.h} &
Defines the base \code{Solver} state: time variables, displacement vectors,
forces, residuals, stiffness matrices, tolerances, energies, and the common
\code{solveStep()} wrapper. \\

\path{src/solver/src/solver.cpp} &
Implements base initialization, force computation, energy accumulation,
step preparation, step commit, perturbations, and termination handling. \\

\path{src/solver/src/solver_implicit.h} &
Declares \code{SteadyStateLinearSolver} and
\code{SteadyStateNonLinearSolver}.  This header also contains the nonlinear
control parameters, default tolerances, matrix-update settings, adaptive
linear-solve settings, and indirect displacement control state. \\

\path{src/solver/src/solver_implicit.cpp} &
Implements the nonlinear solve loop, convergence checks, restart logic,
matrix-update policy, effective matrix construction, linear solver creation,
profiling wrappers, reset path, and adaptive linear tolerance logic. \\

\path{src/solver/src/node_container.cpp} &
Implements full-to-reduced DoF mapping, prescribed boundary updates, constraint
handling, reduced force extraction, external reaction update, and simplex
volumetric strain update. \\

\path{src/solver/src/element_container.cpp} &
Implements internal force integration, element strain/stress evaluation,
material status updates, structural matrix pattern preparation, and stiffness
matrix assembly. \\

\path{src/solver/src/element.cpp} and \path{src/solver/src/element.h} &
Define element-level strain, stress, internal force, and stiffness
contributions.  Specialized element types can override parts of this behavior. \\

\path{src/solver/src/linalg.cpp} and
\path{src/solver/src/linalg_profile.*} &
Define the direct and iterative linear solver wrappers, including profiling
events used by the nonlinear loop. \\
\bottomrule
\end{longtable}

The discussion below focuses on \code{SteadyStateNonLinearSolver}.  Transient
mechanical solvers inherit the same general structure but override parts of the
force and energy evaluation to include damping, inertia, and kinetic energy.

\section{High-Level Execution Model}

At the top level, \oas{} advances through accepted load/time steps.  Each step
is solved by a nonlinear iteration loop.  In the ordinary direct-control path,
one linearized correction is computed per nonlinear iteration:
\[
  K_{\mathrm{eff}} \, \Delta r_f = R_f ,
\]
where \(R_f\) is the reduced residual on free degrees of freedom and
\(\Delta r_f\) is the reduced correction.  The correction is expanded into the
full degree-of-freedom space, constraints are applied, the trial unknown vector
is updated, new internal and external forces are evaluated, and convergence
errors are checked.

The model-level loop is:

\begin{Verbatim}[fontsize=\small, frame=single]
export initial state
while solver is not terminated:
    solver.solveStep()
    export accepted step
    print step timing
\end{Verbatim}

The solver-level \code{solveStep()} wrapper is common to all solvers:

\begin{Verbatim}[fontsize=\small, frame=single]
runBeforeEachStep()
solve()
runAfterEachStep()
\end{Verbatim}

For the nonlinear solver, \code{solve()} may restart the same step with a
smaller time increment if convergence fails.  Only after \code{solve()} returns
does \code{runAfterEachStep()} commit the accepted state.

\section{Notation and Solver Variables}

\subsection{Global Time and Counters}

\begin{longtable}{p{0.20\linewidth} p{0.24\linewidth} p{0.47\linewidth}}
\caption{Time-step and iteration notation.}
\label{tab:time-notation}\\
\toprule
Symbol & C++ member or local & Meaning \\
\midrule
\endfirsthead
\toprule
Symbol & C++ member or local & Meaning \\
\midrule
\endhead
\(n\) & \code{step} &
Accepted load/time-step index.  It is incremented in
\code{Solver::runBeforeEachStep()}. \\

\(t_n\) & \code{time} &
Target time/load coordinate of the current step.  The solver may snap this to
important boundary-condition or function extrema. \\

\(\Delta t_n\) & \code{dt} &
Current time/load increment.  It can be enlarged, reduced, or critically
reduced after failed convergence. \\

\(\Delta t_{\min}\), \(\Delta t_{\max}\) &
\code{dtmin}, \code{dtmax} &
Lower and upper bounds used by adaptive time stepping. \\

\(i\) & \code{it} &
Nonlinear iteration counter inside the current step.  It starts at zero after
each restart. \\

\(i_{\mathrm{cum}}\) & \code{cumul_it} &
Cumulative nonlinear iteration counter within the current call to
\code{solve()}.  It does not reset after a normal iteration, and it is used by
the cumulative stiffness-update rule. \\

\(N_{\max}\) & \code{maxIt} &
Maximum allowed nonlinear iterations before restart or failure. \\

\(N_{\min}\) & \code{minIt} &
Minimum nonlinear iterations required before a converged step can be accepted. \\

\(N_{\mathrm{enlarge}}\) & \code{enlargeIt} &
If convergence is reached before this iteration count, the next step may be
enlarged. \\

\(N_{\mathrm{shorten}}\) & \code{shortenIt} &
If convergence is reached after this iteration count, the next step may be
shortened. \\
\bottomrule
\end{longtable}

\subsection{Unknowns, Corrections, and Forces}

The code distinguishes full vectors, which contain all nodal degrees of freedom,
from reduced vectors, which contain only free degrees of freedom.  Constrained
or prescribed degrees of freedom are recovered through the node and constraint
containers.

\begin{longtable}{p{0.20\linewidth} p{0.24\linewidth} p{0.47\linewidth}}
\caption{Primary nonlinear vectors.}
\label{tab:vector-notation}\\
\toprule
Symbol & C++ member & Meaning \\
\midrule
\endfirsthead
\toprule
Symbol & C++ member & Meaning \\
\midrule
\endhead
\(r_n\) & \code{r} &
Committed full solution vector from the last accepted step.  It is updated only
after the current step is accepted. \\

\(r^{\mathrm{tr}}\) & \code{trial_r} &
Current trial full solution vector inside the nonlinear loop. \\

\(\Delta r_f\) & \code{ddr} &
Reduced correction on free degrees of freedom, returned by the linear solver. \\

\(\Delta r\) & \code{full_ddr} &
Full-space correction obtained by expanding \code{ddr} and applying dependency
relations. \\

\(f_{\mathrm{int}}\) & \code{f_int} &
Full internal force vector assembled from element stresses and internal
sources. \\

\(f_{\mathrm{ext}}\) & \code{f_ext} &
Full external force vector after nodal loads, prescribed-reaction updates, and
constraint reaction contributions are applied. \\

\(f_{\mathrm{load}}\) & \code{load} &
Full prescribed load vector assembled from nodal and body loads before reaction
completion. \\

\(f_{\mathrm{dam}}\) & \code{f_dam} &
Damping force vector.  In the steady-state nonlinear path this remains zero;
transient solvers can add damping contributions. \\

\(f_{\mathrm{acc}}\) & \code{f_acc} &
Inertial force vector.  In the steady-state nonlinear path this remains zero;
transient solvers can add acceleration contributions. \\

\(R\) & \code{residuals} &
Full residual vector.  In the steady-state nonlinear path,
\(R = f_{\mathrm{ext}} - f_{\mathrm{int}}\). \\

\(R_f\) & \code{f} &
Reduced right-hand side extracted from \code{residuals} for free degrees of
freedom after constraint master-force transfer. \\

\(f_{\mathrm{int},n}\), \(f_{\mathrm{ext},n}\) &
\code{f_int_old}, \code{f_ext_old} &
Committed force vectors from the last accepted step.  They are used for energy
integration and restart restoration. \\
\bottomrule
\end{longtable}

\subsection{Matrices}

\begin{longtable}{p{0.20\linewidth} p{0.24\linewidth} p{0.47\linewidth}}
\caption{Matrix notation.}
\label{tab:matrix-notation}\\
\toprule
Symbol & C++ member & Meaning \\
\midrule
\endfirsthead
\toprule
Symbol & C++ member & Meaning \\
\midrule
\endhead
\(K\) & \code{K} &
Structural tangent, secant, elastic, or consistent matrix assembled from
elements before the final effective-system transformation. \\

\(K_{\mathrm{eff}}\) & \code{Keff} &
Effective matrix solved by the linear solver.  In the steady-state nonlinear
path it starts as \(K\), then constraints are transformed into the active
system space.  Transient solvers can add damping or mass terms. \\

\(K_e\) & element local matrix &
Element contribution returned by \code{Element::giveStiffnessMatrix()} or
related damping/mass methods. \\

\(P\) & implicit node/constraint map &
The full-to-reduced and reduced-to-full mapping implemented by
\code{NodeContainer} and \code{ConstraintContainer}.  The code does not store a
single dense matrix \(P\); it applies the map procedurally. \\
\bottomrule
\end{longtable}

\section{Initialization}

\subsection{Model Initialization}

\code{Model::init()} prepares the simulation in a fixed order:

\begin{enumerate}[leftmargin=2em]
  \item Materials are initialized.
  \item Elements are initialized and preprocessing blocks are run.
  \item Boundary-condition objects are initialized.
  \item Nodes are initialized.
  \item Simplex connectivity is initialized when the material setup requires
        simplex volumetric strain support.
  \item Constraints and boundary conditions establish their active degree-of-
        freedom relations.
  \item Fiber/material auxiliary assignments are applied.
  \item Element neighbor or friend relations are initialized.
  \item The solver is initialized.
  \item Exporters are initialized.
\end{enumerate}

This order matters because the solver initialization needs a finalized degree-
of-freedom numbering, element list, material statuses, and boundary-condition
state.

\subsection{Base Solver Initialization}

\code{Solver::init()} allocates all primary full and reduced vectors:
\[
  r, r^{\mathrm{tr}}, f_{\mathrm{int}}, f_{\mathrm{ext}}, f_{\mathrm{dam}},
  f_{\mathrm{acc}}, R \in \R^{N_{\mathrm{dof}}},
\]
and
\[
  \Delta r_f, R_f \in \R^{N_{\mathrm{free}}}.
\]
Initial conditions are read if an initial-condition file is configured;
otherwise the state is set to zero.  The trial vector is initialized from the
committed vector:
\[
  r^{\mathrm{tr}} \leftarrow r.
\]
The solver then computes initial internal and external forces:
\[
  \code{computeInternalExternalForces}(r, f_{\mathrm{load}}, \texttt{frozen=false}, 0, -1).
\]
Finally, element material statuses are committed through
\code{elems->updateMaterialStatuses()}.

\subsection{Linear and Nonlinear Solver Initialization}

\code{SteadyStateLinearSolver::init()} extends the base initialization by:

\begin{enumerate}[leftmargin=2em]
  \item Initializing runtime and linear-solver profilers.
  \item Preparing the sparsity pattern through
        \code{elems->prepareStiffnessMatrix(K)}.
  \item Calling \code{updateSystemMatrices(0, 0, true)} to assemble the first
        structural matrix.
  \item Calling \code{computeKeff()} to build, transform, and factorize the
        first effective system matrix.
\end{enumerate}

\code{SteadyStateNonLinearSolver::init()} additionally:

\begin{enumerate}[leftmargin=2em]
  \item Loads material-status restart data if requested.
  \item Initializes nonlinear convergence state.
  \item Initializes indirect displacement control if configured.
  \item Evaluates forces at the initial integration time with frozen variables
        to initialize field quantities consistently.
\end{enumerate}

\section{Step Preparation}

Every accepted-step attempt starts with \code{runBeforeEachStep()}.  The base
implementation performs the following operations:

\begin{enumerate}[leftmargin=2em]
  \item Increment the step counter:
        \[
          n \leftarrow n + 1.
        \]
  \item Choose the next target time with \code{setNextStepTime()}.
  \item Store the old load vector:
        \[
          f_{\mathrm{load},n} \leftarrow f_{\mathrm{load}}.
        \]
  \item Reset the current load vector and reduced correction:
        \[
          f_{\mathrm{load}} \leftarrow 0,\qquad
          \Delta r_f \leftarrow 0.
        \]
\end{enumerate}

\code{setNextStepTime()} normally advances \(t\) by \(\Delta t\), but it also
looks ahead to the next important time from time-dependent functions and
boundary conditions.  If such a point lies within \(t + 1.25 \Delta t\), the
step target snaps to that exact point.  If the point is a boundary-condition
stage transition, the model jumps to the next boundary-condition stage.

\code{SteadyStateLinearSolver::runBeforeEachStep()} sets:
\[
  r^{\mathrm{tr}} \leftarrow r,
\]
and prints the step header.  \code{SteadyStateNonLinearSolver::runBeforeEachStep()}
adds initialization for indirect displacement control and prints the nonlinear
error tolerance header.

\section{Nonlinear Solver Algorithm}

\subsection{Direct-Control Path}

The ordinary nonlinear path is used when no indirect displacement control object
is configured.  The structure of \code{SteadyStateNonLinearSolver::solve()} is:

\begin{Verbatim}[fontsize=\small, frame=single]
converged = false
restarted = false
while not converged:
    assemble load and boundary conditions for current time
    compute frozen forces at integration time
    it = 0
    while not converged and it < maxIt:
        maybe update stiffness matrix and factorize Keff
        reduce full residual R to reduced RHS R_f
        solve Keff * ddr = R_f
        expand ddr to full_ddr and update trial_r
        compute unfrozen forces at updated trial_r
        evaluate residual, displacement, and energy errors
        decide whether convergence is acceptable
        export iteration data
        it += 1
    if converged:
        compute step-end forces
    else:
        restart, relax tolerance, or terminate
    adapt next dt if accepted
\end{Verbatim}

The nonlinear iteration is Newton-like: it repeatedly solves a linearized
problem, updates the trial state, and checks whether residual, increment, and
energy measures are small enough.  The exact meaning of the matrix depends on
the configured stiffness matrix type and update frequency.  With a tangent
matrix and frequent updates the method is closest to Newton's method.  With
secant or less frequent updates it becomes a modified Newton or quasi-Newton
procedure.

\subsection{Iteration Start: Matrix Update}

At the beginning of every nonlinear iteration, the solver may update the
structural matrix:
\[
  K \leftarrow \sum_e K_e,
\]
then build and factorize
\[
  K_{\mathrm{eff}}.
\]
The update is controlled by \code{updateSystemMatrices(iteration, cumul_iteration, enforce)}.
The matrix is updated if at least one of the following conditions is true:

\begin{itemize}[leftmargin=2em]
  \item \code{enforce} is true.
  \item The current step/iteration satisfies the configured step-update rule.
  \item \code{stiffnessMatrixIterUpdate == 0}, meaning every iteration.
  \item A positive iteration-period rule matches the current iteration.
  \item A positive cumulative-iteration-period rule matches the cumulative
        counter.
\end{itemize}

When the matrix is updated, the first nonlinear iteration can use a special
first-iteration matrix type if \code{stiffMatTypeFirstIT} is not \code{void}.
Otherwise it uses \code{stiffMatType}.  Valid matrix types are shown in
Table~\ref{tab:matrix-types}.

\begin{longtable}{p{0.20\linewidth} p{0.72\linewidth}}
\caption{Configured stiffness matrix types.}
\label{tab:matrix-types}\\
\toprule
Type & Meaning in the code \\
\midrule
\endfirsthead
\toprule
Type & Meaning in the code \\
\midrule
\endhead
\code{elastic} &
Uses an elastic stiffness contribution from material statuses.  This is usually
more robust but less nonlinear-state-specific. \\

\code{secant} &
Uses a secant material stiffness.  This is the default in the nonlinear solver. \\

\code{tangent} &
Uses a tangent stiffness, when supplied by the material model.  This is the
closest option to a classical Newton tangent. \\

\code{consistent} &
Uses a consistent material stiffness contribution, when supplied. \\
\bottomrule
\end{longtable}

After assembly, \code{computeKeff()} copies:
\[
  K_{\mathrm{eff}} \leftarrow K,
\]
then transforms the matrix into constraint space if constraints are active.
Finally, \code{factorizeLinearSystem()} analyzes or factorizes the selected
linear solver backend.

\subsection{Residual Reduction}

The force vector used by the linear solver is not the raw full residual.  The
solver calls:
\[
  \code{nodes->giveReducedForceArray(residuals, f)}.
\]
This first transfers slave or constrained forces back to their master degrees of
freedom:
\[
  R \leftarrow \text{constraint-master-force-transform}(R),
\]
then extracts only the free active degrees of freedom:
\[
  R_f = R|_{\mathrm{free}}.
\]
This reduced vector is the right-hand side of the linear solve.

\subsection{Linear Correction}

The direct-control correction is:
\[
  K_{\mathrm{eff}} \Delta r_f = R_f .
\]
In code this is:
\[
  \code{profiledLinearSolve(ddr, f, "direct_residual")}.
\]
The wrapper records timing and linear-solver metadata.  For iterative solvers it
also records iteration counts, requested tolerance, achieved error, deflation
statistics, and specialized sub-phase timings if available.

If the solve fails, the nonlinear step is terminated immediately.  If the solve
succeeds and the solver supports deflation vectors, the correction from a
successful direct residual solve can be collected:
\[
  \code{collectLinearDeflationVector(ddr, "direct_residual", success)}.
\]

\subsection{Trial-State Update}

The reduced correction is expanded into the full degree-of-freedom space:
\[
  \Delta r \leftarrow P \Delta r_f ,
\]
through:
\[
  \code{nodes->giveFullDoFArray(ddr, full_ddr)}.
\]
The node and constraint layer then accounts for dependency relations that depend
on conjugate variables:
\[
  \code{nodes->updateFullDoFsByDependenciesOnConjugates(full_ddr, trial_r, f_ext)}.
\]
Finally:
\[
  r^{\mathrm{tr}} \leftarrow r^{\mathrm{tr}} + \Delta r.
\]

\subsection{Force Re-Evaluation}

After the trial vector is updated, the solver recomputes forces at the new
trial state:
\[
  \code{computeForcesAtIntegrationTime(false)}.
\]
For the steady-state nonlinear solver this calls:
\[
  \code{computeInternalExternalForces(trial_r, load, frozen=false, time, it)}.
\]
The resulting steady residual is:
\[
  R = f_{\mathrm{ext}} - f_{\mathrm{int}}.
\]

\section{Force Assembly and Residual Construction}

\subsection{Internal and External Force Routine}

\code{Solver::computeInternalExternalForces()} performs three main operations:

\begin{enumerate}[leftmargin=2em]
  \item Update simplex volumetric strains:
        \[
          \code{nodes->updateSimplexVolumetricStrains(rr)}.
        \]
  \item Integrate element internal forces:
        \[
          \code{elems->integrateInternalForces(rr, f_int, frozen, t, timeStep)}.
        \]
  \item Complete external forces and reactions:
        \[
          \code{nodes->updateExternalForcesByReactions(...)}.
        \]
\end{enumerate}

The steady residual is then formed explicitly:
\[
  R = f_{\mathrm{ext}} - f_{\mathrm{int}}.
\]

\subsection{Element Internal Force Integration}

\code{ElementContainer::integrateInternalForces()} assembles the full internal
force vector.  Its conceptual algorithm is:

\begin{Verbatim}[fontsize=\small, frame=single]
f_int = 0
reset eigenstrain contributions
run material preparation for stress evaluation
for solution_order = 0 ... max_solution_order:
    for each element with this solution_order:
        gather element dof values from trial_r
        evaluate element strains
        evaluate element stresses
for each element:
    compute element internal force vector
    scatter local force entries into global f_int
\end{Verbatim}

The \code{solution_order} loop is important for coupled or dependent element
quantities.  It enforces a deterministic order between element groups whose
stress/strain evaluation may depend on previously updated fields.

The \code{frozen} argument controls material-state handling:

\begin{itemize}[leftmargin=2em]
  \item \code{frozen=true}: stresses are evaluated with frozen internal
        variables.  This is used for predictor-style or initialization passes
        where the code wants forces without advancing temporary material state.
  \item \code{frozen=false}: stresses are evaluated normally for the current
        trial state, updating temporary material-status quantities used by the
        current nonlinear iteration.
\end{itemize}

\subsection{External Force and Reaction Completion}

The initial \code{load} vector contains applied nodal and body-load terms.  It
is not yet the final \code{f_ext}.  The node container completes the external
force vector by:

\begin{enumerate}[leftmargin=2em]
  \item Transforming internal, damping, acceleration, and load vectors through
        constraint master-force relations.
  \item Starting from:
        \[
          f_{\mathrm{ext}} \leftarrow f_{\mathrm{load}}.
        \]
  \item Adding Lagrange-multiplier reaction terms to master degrees of freedom.
  \item For directly prescribed blocked degrees of freedom, setting the
        external force equal to the force required for equilibrium:
        \[
          f_{\mathrm{ext},i}
          =
          f_{\mathrm{int},i}
          +
          f_{\mathrm{dam},i}
          +
          f_{\mathrm{acc},i}.
        \]
\end{enumerate}

This is why the full residual should vanish on prescribed degrees of freedom:
their reaction force is defined from the internal, damping, and inertial force
balance.

\section{Boundary Conditions, Constraints, and DoF Mapping}

\subsection{Degree-of-Freedom Numbering}

\code{NodeContainer::establishDoFArray()} builds the solver numbering.  The
layout is:

\begin{enumerate}[leftmargin=2em]
  \item Free degrees of freedom.
  \item Slave constrained degrees of freedom.
  \item Directly blocked or prescribed degrees of freedom.
\end{enumerate}

Only the first group is part of the reduced linear system.  The inverse map
\code{invDoFid} maps reduced indices back to full degree-of-freedom indices.
The physical field of every degree of freedom is stored in
\code{physicalFieldsDoF}; this is later used by the convergence-error
calculation.

\subsection{Boundary Updates}

For direct control, at the start of an attempted step:
\[
  \code{nodes->addRHS_nodalLoad(load, time)}
\]
adds active body and nodal loads at the current time.  Then:
\[
  \code{nodes->updateDirrichletBC(trial_r, time)}
\]
writes prescribed displacement or field values into the full trial vector and
updates dependent constrained degrees of freedom.

\subsection{Reduced and Full-Space Operations}

The most important map operations are:

\begin{longtable}{p{0.48\linewidth} p{0.44\linewidth}}
\caption{Node and constraint mapping routines used by the nonlinear solver.}
\label{tab:mapping-routines}\\
\toprule
Routine & Meaning \\
\midrule
\endfirsthead
\toprule
Routine & Meaning \\
\midrule
\endhead
\code{giveReducedForceArray(fullf, f)} &
Transfers constrained slave forces to masters, then extracts free rows into a
reduced vector. \\

\code{giveFullDoFArray(fDoFs, fullDoFs)} &
Expands a reduced correction into the full vector and calculates dependent
constrained degrees of freedom. \\

\code{updateFullDoFsByDependenciesOnConjugates(...)} &
Applies additional correction terms for constraints whose dependent values are
functions of conjugate variables such as forces. \\

\code{updateExternalForcesByReactions(...)} &
Completes \code{f_ext} by adding reactions and making prescribed degrees of
freedom force-balanced. \\
\bottomrule
\end{longtable}

\section{Convergence Criteria}

\subsection{Physical-Field Separation}

The nonlinear solver computes convergence errors separately for each physical
field, then combines them.  The current field indices are listed in
Table~\ref{tab:physical-fields}.

\begin{longtable}{p{0.16\linewidth} p{0.22\linewidth} p{0.46\linewidth}}
\caption{Physical fields used by the error calculation.}
\label{tab:physical-fields}\\
\toprule
Index & Field & Stabilizing floor \(\epsilon_p^2\) \\
\midrule
\endfirsthead
\toprule
Index & Field & Stabilizing floor \(\epsilon_p^2\) \\
\midrule
\endhead
0 & Mechanics & \(10^{-20}\) \\
1 & Transport & \(10^{-10}\) \\
2 & Temperature & \(10^{-17}\) \\
3 & Humidity & \(10^{-18}\) \\
\bottomrule
\end{longtable}

Let \(I_p\) be the set of full degrees of freedom belonging to physical field
\(p\).  The code accumulates squared norms:
\[
  S_R^p = \sum_{i\in I_p} R_i^2,
  \qquad
  S_{\Delta r}^p = \sum_{i\in I_p} (\Delta r_i)^2,
  \qquad
  S_r^p = \sum_{i\in I_p} (r_i^{\mathrm{tr}})^2.
\]
It also accumulates force scales:
\[
  S_{\mathrm{ext}}^p = \sum_{i\in I_p} f_{\mathrm{ext},i}^2,\quad
  S_{\mathrm{int}}^p = \sum_{i\in I_p} f_{\mathrm{int},i}^2,\quad
  S_{\mathrm{dam}}^p = \sum_{i\in I_p} f_{\mathrm{dam},i}^2,\quad
  S_{\mathrm{acc}}^p = \sum_{i\in I_p} f_{\mathrm{acc},i}^2.
\]
For energy, it accumulates:
\[
  E_R^p = \sum_{i\in I_p} R_i \Delta r_i.
\]

Some elements or RVE/eigenstrain mechanisms can add external error
contributions to the displacement, force, and external-work scales.  The code
adds these before final normalization.

\subsection{Residual Error}

The residual error is:
\[
  e_R =
  \sqrt{
  \sum_p
  \frac{S_R^p}
       {\max\left(
          S_{\mathrm{ext}}^p,
          S_{\mathrm{int}}^p,
          S_{\mathrm{dam}}^p,
          S_{\mathrm{acc}}^p,
          \epsilon_p^2
       \right)}
  }.
\]
In code this value is stored in \code{resErr}.  It measures force imbalance
relative to the largest relevant force scale in each physical field.

\subsection{Displacement or Increment Error}

The displacement increment error is:
\[
  e_{\Delta r} =
  \sqrt{
  \sum_p
  \frac{S_{\Delta r}^p}
       {\max\left(S_r^p, \epsilon_p^2\right)}
  }.
\]
In code this value is stored in \code{disErr}.  On the first nonlinear
iteration of a step, the code sets the displayed and tested displacement error
to zero after evaluation:
\[
  i = 0 \quad \Rightarrow \quad e_{\Delta r} \leftarrow 0.
\]
This avoids rejecting the first correction based only on an increment compared
with a still-developing trial scale.

\subsection{Energy Error}

The energy error is:
\[
  e_E =
  \sqrt{
  \sum_p
  \frac{\abs{E_R^p}}
       {\max\left(
          \abs{W_{\mathrm{ext}}^p},
          \abs{W_{\mathrm{int}}^p},
          \abs{W_{\mathrm{kin}}^p},
          \epsilon_p^2
       \right)}
  }.
\]
In code this value is stored in \code{eneErr}.  It measures residual work during
the latest correction relative to accumulated external, internal, or kinetic
energy scales.

\subsection{Convergence Decision}

The basic convergence test is:
\[
  e_{\Delta r} \leq \tau_{\Delta r},
  \qquad
  e_R \leq \tau_R,
  \qquad
  e_E \leq \tau_E,
  \qquad
  i+1 \geq N_{\min}.
\]
The tolerances are:

\begin{longtable}{p{0.22\linewidth} p{0.25\linewidth} p{0.44\linewidth}}
\caption{Nonlinear tolerances.}
\label{tab:tolerances}\\
\toprule
Quantity & C++ member & Default \\
\midrule
\endfirsthead
\toprule
Quantity & C++ member & Default \\
\midrule
\endhead
\(\tau_{\Delta r}\) & \code{maxDisErr} & \(10^{-5}\) \\
\(\tau_R\) & \code{maxResErr} & \(10^{-5}\) \\
\(\tau_E\) & \code{maxEneErr} & \(10^{-5}\) \\
\(\tau_{\Delta r}^{\mathrm{limit}}\) & \code{limitDisErr} & 0 unless configured \\
\(\tau_R^{\mathrm{limit}}\) & \code{limitResErr} & 0 unless configured \\
\(\tau_E^{\mathrm{limit}}\) & \code{limitEneErr} & 0 unless configured \\
\bottomrule
\end{longtable}

If adaptive linear tolerance is active, the nonlinear solver can require one
additional correction with a tighter linear tolerance before accepting a step
that otherwise satisfies the nonlinear errors.  This prevents a loose linear
solve from being the last correction in a nominally converged nonlinear step
when strict final convergence is requested.

If any of the nonlinear errors is NaN, the current attempt is treated as failed.
The iteration counter is forced to the maximum and a large residual error is
assigned.

\section{Adaptive Linear Tolerance}

The nonlinear solver can dynamically switch the requested tolerance of the
linear solver during direct residual solves.  This is controlled by the
adaptive-linear-tolerance members, especially:

\begin{itemize}[leftmargin=2em]
  \item \code{linearAdaptiveTolEnabled};
  \item \code{linearAdaptiveLooseTol};
  \item \code{linearAdaptiveTightTol};
  \item \code{linearAdaptiveTriggerRatio};
  \item \code{linearAdaptiveRequireTightConvergence}.
\end{itemize}

The rule is:

\begin{enumerate}[leftmargin=2em]
  \item Adaptive tolerance only applies to RHS kind \code{direct_residual}.
  \item The loose tolerance is used by default.
  \item The tight tolerance is used if the solver has requested a strict final
        solve or if the current nonlinear error ratio is already below the
        configured trigger ratio.
  \item The nonlinear error ratio is the maximum of the residual, energy, and
        displacement errors divided by their target tolerances.  The
        displacement component is ignored for the first nonlinear iteration.
  \item If a step appears converged after a loose linear solve and strict final
        convergence is required, the solver performs one more iteration using
        the tight linear tolerance.
\end{enumerate}

This mechanism is meant to avoid oversolving early nonlinear iterations while
still ending each accepted nonlinear step with a sufficiently accurate linear
correction.

\section{Indirect Displacement Control}

When \code{idc} is configured, the nonlinear solver does not simply solve one
linear correction per iteration.  It solves two related systems and computes a
load multiplier correction.

The variables are:

\begin{longtable}{p{0.22\linewidth} p{0.28\linewidth} p{0.41\linewidth}}
\caption{Indirect displacement control variables.}
\label{tab:idc-vars}\\
\toprule
Symbol & C++ member & Meaning \\
\midrule
\endfirsthead
\toprule
Symbol & C++ member & Meaning \\
\midrule
\endhead
\(\lambda\) & \code{idc_time} &
The indirect-control load coordinate currently used to evaluate loads and
boundary conditions. \\

\(\lambda_n\) & \code{idc_time_converged} &
Last accepted indirect-control load coordinate. \\

\(\Delta \lambda_0\) & \code{idc_dt} &
Small probe increment used to form a load-increment direction. \\

\(\mu\) & \code{load_mult} &
Multiplier correction returned by the indirect-control object. \\

\(\Delta r_f^R\) & \code{ddr} &
Residual correction direction. \\

\(\Delta r_f^L\) & \code{ddf} &
Load-increment correction direction. \\
\bottomrule
\end{longtable}

The conceptual sequence is:

\begin{enumerate}[leftmargin=2em]
  \item Save the current trial state and reduced residual.
  \item Solve:
        \[
          K_{\mathrm{eff}} \Delta r_f^R = R_f.
        \]
  \item Advance the indirect-control load coordinate by \(\Delta \lambda_0\),
        update loads and boundary conditions, and recompute forces.
  \item Solve:
        \[
          K_{\mathrm{eff}} \Delta r_f^L
          =
          R_f(\lambda+\Delta\lambda_0) - R_f(\lambda).
        \]
  \item Form trial states A and B, optionally evaluate external forces for
        both states, and ask the indirect-control object for the multiplier
        \(\mu\).
  \item Correct the residual direction:
        \[
          \Delta r_f \leftarrow \Delta r_f^R + \mu \Delta r_f^L.
        \]
  \item Update:
        \[
          \lambda \leftarrow \lambda + \mu \Delta\lambda_0.
        \]
  \item Restore the original trial state, apply the corrected load and boundary
        conditions at the new indirect-control coordinate, then apply the
        corrected displacement increment.
\end{enumerate}

On restart, \code{idc_time} is reset to \code{idc_time_converged}, so the failed
trial load advancement is discarded.

\section{Restart, Relaxed Acceptance, and Termination}

When the inner nonlinear loop ends, the solver has either converged or exhausted
the nonlinear iteration budget.

\subsection{Normal Convergence}

If the convergence criteria are satisfied, the solver marks:
\[
  \code{fully_converged} \leftarrow \code{true},
\]
and computes step-end forces:
\[
  \code{computeForcesAtStepEnd(false)}.
\]
For the steady-state nonlinear solver, this is another force evaluation at the
accepted trial state.  It ensures the final force state is consistent before the
step is committed.

\subsection{Critical Restart}

If the step did not converge and the current increment is larger than the
minimum increment, the solver restarts the same step:
\[
  \Delta t \leftarrow
  \max(\gamma_{\mathrm{crit}} \Delta t, \Delta t_{\min}),
\]
where:
\[
  \gamma_{\mathrm{crit}} = \code{critical_step_decrease}.
\]
The default is \(0.5\).

The restart restores:

\begin{itemize}[leftmargin=2em]
  \item the target time to the previous value, then advances it by the reduced
        \(\Delta t\);
  \item \code{trial_r} back to the last committed \code{r};
  \item old internal and external forces;
  \item the load vector and reduced correction to zero;
  \item material statuses through \code{elems->resetMaterialStatuses()};
  \item indirect-control state, if active.
\end{itemize}

The nonlinear loop then tries the same physical step again with the smaller
increment.

\subsection{Relaxed-Tolerance Acceptance}

If the step still did not meet the main tolerances, but the errors satisfy the
configured limit tolerances, the solver can accept the step as not fully
converged:
\[
  \code{fully_converged} \leftarrow \code{false}.
\]
This path is only available when nonzero limit tolerances were configured.
Otherwise, the default limit tolerances are zero and this fallback is disabled.

\subsection{Failure}

If the step cannot be restarted and cannot be accepted under relaxed tolerances,
the solver prints an error message and sets:
\[
  \code{terminated} \leftarrow \code{true}.
\]
The model-level solve loop then stops.

\section{Adaptive Time-Step Control}

After an accepted step, the next increment can be adjusted based on the number
of nonlinear iterations:

\[
\Delta t_{\mathrm{next}} =
\begin{cases}
\min(\gamma_{\mathrm{inc}} \Delta t, \Delta t_{\max}),
  & \text{if } i < N_{\mathrm{enlarge}} \text{ and no restart occurred},\\
\max(\gamma_{\mathrm{dec}} \Delta t, \Delta t_{\min}),
  & \text{if } i > N_{\mathrm{shorten}},\\
\Delta t,
  & \text{otherwise}.
\end{cases}
\]

The default multipliers are:

\begin{longtable}{p{0.27\linewidth} p{0.27\linewidth} p{0.35\linewidth}}
\caption{Adaptive time-step factors.}
\label{tab:timestep-factors}\\
\toprule
Meaning & C++ member & Default \\
\midrule
\endfirsthead
\toprule
Meaning & C++ member & Default \\
\midrule
\endhead
Normal enlargement factor & \code{step_increase} & \(1.25\) \\
Normal reduction factor & \code{step_decrease} & \(0.8\) \\
Critical restart reduction factor & \code{critical_step_decrease} & \(0.5\) \\
\bottomrule
\end{longtable}

The final increment is clamped to \([\Delta t_{\min}, \Delta t_{\max}]\).  The
base \code{runAfterEachStep()} also reduces the next increment if it would pass
the configured termination time.

\section{Step Commit}

After \code{solve()} returns, \code{solveStep()} calls
\code{runAfterEachStep()}.  The base implementation performs the actual state
commit:

\begin{enumerate}[leftmargin=2em]
  \item Apply scheduled perturbations if they are active at the current time.
  \item Compute total internal, external, and kinetic energies.
  \item Commit the trial vector:
        \[
          r \leftarrow r^{\mathrm{tr}}.
        \]
  \item Store old force vectors:
        \[
          f_{\mathrm{int},n} \leftarrow f_{\mathrm{int}},\qquad
          f_{\mathrm{ext},n} \leftarrow f_{\mathrm{ext}}.
        \]
  \item Commit element material statuses:
        \[
          \code{elems->updateMaterialStatuses()}.
        \]
  \item Clamp or terminate based on the next \(\Delta t\).
  \item Commit old energy values.
\end{enumerate}

The nonlinear override prints a separator and, for indirect displacement
control, commits:
\[
  \code{idc_time_converged} \leftarrow \code{idc_time}.
\]

\section{Energy Accounting}

The base energy update integrates internal and external work with a trapezoidal
rule over the accepted correction:
\[
  W_{\mathrm{int}}^p
  \leftarrow
  W_{\mathrm{int},n}^p
  +
  \sum_{i\in I_p}
  \frac{1}{2}
  \left(f_{\mathrm{int},i} + f_{\mathrm{int},n,i}\right)
  \left(r_i^{\mathrm{tr}} - r_{n,i}\right),
\]
\[
  W_{\mathrm{ext}}^p
  \leftarrow
  W_{\mathrm{ext},n}^p
  +
  \sum_{i\in I_p}
  \frac{1}{2}
  \left(f_{\mathrm{ext},i} + f_{\mathrm{ext},n,i}\right)
  \left(r_i^{\mathrm{tr}} - r_{n,i}\right).
\]
In the steady-state base class, kinetic energy is not updated.  Transient
mechanical solvers can override this behavior to include a meaningful
\(W_{\mathrm{kin}}\).

\section{Matrix Assembly and Effective System}

\subsection{Element Stiffness Assembly}

The element container prepares the sparsity pattern before the nonlinear solve
starts.  During a matrix update it:

\begin{enumerate}[leftmargin=2em]
  \item Resets matrix values while preserving the pattern.
  \item Adds Lagrange multiplier constraint entries.
  \item Requests local element matrices \(K_e\) from each element.
  \item Scatters local entries into global matrix entries using element
        degree-of-freedom indices.
\end{enumerate}

The code supports optimized assembly paths using a precomputed sparse-entry map
when the current matrix pattern is compatible.  The mathematical result remains
the same: the global structural matrix is the sum of all element and constraint
contributions.

\subsection{Effective Matrix Construction}

For the steady-state nonlinear solver:
\[
  K_{\mathrm{eff}} \leftarrow K.
\]
If constraints exist:
\[
  K_{\mathrm{eff}} \leftarrow
  \text{constraint-space-transform}(K_{\mathrm{eff}}).
\]
The transformed matrix is then passed to the configured linear solver.  The
same high-level function also handles linear solver creation on first use and
calls symbolic analysis where required.

\section{Linear Solver Interface}

\subsection{Supported Solver Families}

\code{factorizeLinearSystem()} selects the configured solver backend.  The
available choices depend on compile-time flags.  The implemented families
include:

\begin{itemize}[leftmargin=2em]
  \item Eigen direct and iterative solvers.
  \item Pardiso, Cholmod, SuperLU, when compiled in.
  \item AMGCL conjugate-gradient solver, when compiled in.
  \item Hypre BoomerAMG conjugate-gradient solver, when compiled in.
  \item Native \code{DeflatedFGMRES}.
\end{itemize}

For AMGCL, Hypre, and native deflated FGMRES paths, the solver may also build an
elastic degree-of-freedom map used for block structure, nullspace hints, lifting,
or deflation metadata.

\subsection{Profiling Wrapper}

The nonlinear solver does not call the linear solver directly.  It calls:
\[
  \code{profiledLinearSolve(x, b, rhsKind)}.
\]
The wrapper records:

\begin{itemize}[leftmargin=2em]
  \item total linear solve time;
  \item matrix event number and nonlinear iteration context;
  \item right-hand-side kind;
  \item right-hand-side norm and change from the previous RHS;
  \item solution norm;
  \item linear iteration count and reported error;
  \item requested tolerance;
  \item deflation basis size and deflation update metrics, when available;
  \item detailed sub-phase timings for native deflated FGMRES, when available.
\end{itemize}

\begin{longtable}{p{0.30\linewidth} p{0.60\linewidth}}
\caption{RHS kinds used by the nonlinear and reset paths.}
\label{tab:rhs-kinds}\\
\toprule
RHS kind & Meaning \\
\midrule
\endfirsthead
\toprule
RHS kind & Meaning \\
\midrule
\endhead
\code{direct_residual} &
Ordinary nonlinear correction for direct-control steps.  Successful solves can
contribute vectors to the deflation basis. \\

\code{idc_residual} &
Residual correction direction for indirect displacement control. \\

\code{idc_load_increment} &
Load-increment correction direction for indirect displacement control. \\

\code{reset_residual} &
Residual correction used by the nonlinear reset procedure. \\

\code{linear_step} &
Single linear solve used by the steady-state linear solver path. \\
\bottomrule
\end{longtable}

\section{Runtime Profiling Names}

The solver contains nested timing scopes.  The exact set can vary with compile
options and selected solver, but the nonlinear path commonly records the
categories shown in Table~\ref{tab:runtime-phases}.

\begin{longtable}{p{0.34\linewidth} p{0.56\linewidth}}
\caption{Important runtime profile names.}
\label{tab:runtime-phases}\\
\toprule
Profile name & Meaning \\
\midrule
\endfirsthead
\toprule
Profile name & Meaning \\
\midrule
\endhead
\code{model.solve.solve_step} &
Wall time spent in one call to \code{solver->solveStep()}. \\

\code{solver.before_step} &
Step preparation, trial-state reset, time update, and nonlinear before-step
logic. \\

\code{solver.solve} &
Main solver body for the current step. \\

\code{solver.after_step} &
Accepted-step commit, energy update, material status commit, and termination
checks. \\

\code{nonlinear.iteration} &
One nonlinear iteration, including matrix update when triggered, linear solve,
state update, force evaluation, and error evaluation. \\

\code{matrix.stiffness_update} &
Structural stiffness matrix assembly. \\

\code{matrix.compute_keff} &
Effective matrix construction, constraint transform, and factorization. \\

\code{linear.solve_total} &
One profiled linear solve. \\

\code{forces.total} &
Total force evaluation wrapper. \\

\code{forces.simplex_volumetric_strains} &
Simplex volumetric strain update before element force integration. \\

\code{forces.integrate_internal_forces} &
Element strain, stress, and internal-force integration. \\

\code{forces.external_reactions} &
Completion of external forces and reactions after internal force assembly. \\

\code{nonlinear.evaluate_errors} &
Residual, displacement, and energy error calculation. \\
\bottomrule
\end{longtable}

Native deflated FGMRES can further split \code{linear.solve_total} into matvec,
preconditioner setup/apply, orthogonalization, reorthogonalization, deflation
projection, least-squares work, and residual checks.

\section{Reset Path}

\code{SteadyStateNonLinearSolver::reset()} is a separate nonlinear correction
path used when the model state must be re-equilibrated without advancing the
usual load step.  It follows the same basic pattern:

\begin{enumerate}[leftmargin=2em]
  \item Compute forces with frozen internal variables.
  \item Repeatedly update the matrix as needed.
  \item Reduce the residual.
  \item Solve \code{reset_residual}.
  \item Update the trial vector.
  \item Recompute frozen forces.
  \item Evaluate convergence errors.
\end{enumerate}

Unlike the normal \code{solve()} method, the reset routine calls
\code{runAfterEachStep()} itself after it converges.  It is therefore a
self-contained equilibration and commit operation.

\section{Configuration Summary}

\begin{longtable}{p{0.24\linewidth} p{0.34\linewidth} p{0.34\linewidth}}
\caption{Important nonlinear configuration options read by the solver.}
\label{tab:config}\\
\toprule
Input concept & C++ member & Effect \\
\midrule
\endfirsthead
\toprule
Input concept & C++ member & Effect \\
\midrule
\endhead
Initial time step & \code{initdt}, \code{dt} &
Starting load/time increment. \\

Minimum/maximum time step & \code{dtmin}, \code{dtmax} &
Enable and bound adaptive time stepping. \\

Residual tolerance & \code{maxResErr} &
Target residual error. \\

Energy tolerance & \code{maxEneErr} &
Target residual-work error. \\

Increment tolerance & \code{maxDisErr} &
Target correction-size error. \\

Limit tolerances & \code{limitResErr}, \code{limitEneErr}, \code{limitDisErr} &
Optional relaxed acceptance thresholds. \\

Minimum/maximum nonlinear iterations & \code{minIt}, \code{maxIt} &
Lower and upper nonlinear iteration limits. \\

Step enlarge/shorten iteration counts & \code{enlargeIt}, \code{shortenIt} &
Iteration thresholds for next-step adaptation. \\

Step increase/decrease factors & \code{step_increase}, \code{step_decrease} &
Normal accepted-step time-step scaling. \\

Critical decrease factor & \code{critical_step_decrease} &
Failed-step restart scaling. \\

Matrix update by step & \code{stiffnessMatrixStepUpdate} &
Controls whether the first iteration of some steps updates the matrix. \\

Matrix update by iteration & \code{stiffnessMatrixIterUpdate} &
Controls per-step nonlinear iteration update frequency. \\

Matrix update by cumulative iteration & \code{stiffnessMatrixCumulIterUpdate} &
Controls matrix update frequency across restarts or cumulative iterations. \\

Matrix type & \code{stiffMatType} &
Selects \code{elastic}, \code{secant}, \code{tangent}, or \code{consistent}. \\

First-iteration matrix type & \code{stiffMatTypeFirstIT} &
Optional different matrix type for the first nonlinear iteration. \\

Adaptive linear tolerance & \code{linearAdaptiveTolEnabled} and related members &
Switches iterative linear solve tolerance between loose and tight values inside
the nonlinear process. \\
\bottomrule
\end{longtable}

\section{What a Single Direct-Control Iteration Means Physically}

For mechanics, the nonlinear iteration can be read as the following equilibrium
correction:

\begin{enumerate}[leftmargin=2em]
  \item Given the current trial displacement \(u^{\mathrm{tr}}\), compute the
        internal nodal force \(f_{\mathrm{int}}(u^{\mathrm{tr}})\) from element
        stresses.
  \item Compute the external nodal force \(f_{\mathrm{ext}}\), including
        reactions for prescribed degrees of freedom.
  \item The imbalance is
        \[
          R = f_{\mathrm{ext}} - f_{\mathrm{int}}.
        \]
  \item Reduce that imbalance to the active free system \(R_f\).
  \item Solve a linearized equilibrium correction
        \[
          K_{\mathrm{eff}} \Delta u_f = R_f.
        \]
  \item Expand \(\Delta u_f\), update constrained dependent quantities, and set
        \[
          u^{\mathrm{tr}} \leftarrow u^{\mathrm{tr}} + \Delta u.
        \]
  \item Repeat until the remaining force imbalance, correction size, and
        residual work are all small relative to their natural scales.
\end{enumerate}

For transport, temperature, and humidity fields the same algebraic structure is
used.  The physical meaning of \(r\), \(f_{\mathrm{int}}\), and
\(f_{\mathrm{ext}}\) changes with the field, and the convergence calculation
normalizes each field separately before combining them.

\section{Important Implementation Consequences}

\subsection{The Committed State Does Not Move During Iterations}

During nonlinear iterations:
\[
  r \quad \text{remains fixed}, \qquad
  r^{\mathrm{tr}} \quad \text{changes}.
\]
This distinction is important for restarts.  A failed step can restore
\(r^{\mathrm{tr}} \leftarrow r\) without losing the last accepted state.

\subsection{Material Statuses Have Trial and Committed Roles}

Element force evaluation can update temporary material-status quantities during
the trial process.  These are only committed by
\code{elems->updateMaterialStatuses()} after the step is accepted.  If a step is
restarted, \code{elems->resetMaterialStatuses()} discards the failed trial
state.

\subsection{Prescribed Degrees of Freedom Are Balanced by Reactions}

The residual on directly prescribed degrees of freedom is not used as an
unknown equation in the reduced system.  Instead, the external force on those
degrees of freedom is set to the value needed to balance internal, damping, and
inertial forces.  This makes the reaction force an output quantity rather than
an unknown correction.

\subsection{Matrix Update Frequency Changes the Nonlinear Method}

If the matrix is rebuilt every nonlinear iteration with a tangent stiffness, the
algorithm behaves most like Newton's method.  If the matrix is rebuilt less
often or uses a secant/elastic matrix, the method trades convergence rate for
cheaper iterations and sometimes improved robustness.

\subsection{Linear Solver Accuracy Can Affect Nonlinear Acceptance}

For iterative linear solvers, the correction is approximate.  Adaptive linear
tolerance allows early nonlinear iterations to be solved loosely, but the code
can require a tight final correction before accepting the nonlinear step.  This
is especially relevant for deflated FGMRES and AMG preconditioned paths.

\section{Compact Algorithm Reference}

\begin{Verbatim}[fontsize=\small, frame=single]
Model::solve:
    export initial data
    while not solver.terminated:
        solver.solveStep()
        export accepted step

Solver::solveStep:
    runBeforeEachStep()
    solve()
    runAfterEachStep()

SteadyStateNonLinearSolver::solve, direct control:
    converged = false
    while not converged:
        load = loads(time)
        trial_r = prescribed_boundary_values(trial_r, time)
        updateFieldVariables() with zero ddr
        computeForcesAtIntegrationTime(frozen=true)

        it = 0
        while not converged and it < maxIt:
            if matrix_update_required:
                K = assemble_structural_matrix(stiffness_type)
                Keff = constraint_transform(K)
                factorize_or_setup_linear_solver(Keff)

            f = reduce(residuals)
            ddr = solve(Keff, f)
            full_ddr = expand_and_apply_dependencies(ddr)
            trial_r += full_ddr

            computeForcesAtIntegrationTime(frozen=false)
            evaluateErrors()
            converged = errors_below_tolerances and it+1 >= minIt
            export iteration
            it += 1

        if converged:
            fully_converged = true
            computeForcesAtStepEnd(frozen=false)
        else if dt can be reduced:
            restore committed state
            reset material statuses
            reduce dt and retry
        else if limit tolerances are satisfied:
            fully_converged = false
            computeForcesAtStepEnd(frozen=false)
            converged = true
        else:
            terminated = true
            return

        adapt next dt

Solver::runAfterEachStep:
    compute energies
    r = trial_r
    f_int_old = f_int
    f_ext_old = f_ext
    commit material statuses
    update termination and next dt
\end{Verbatim}

\section{Reading Solver Output}

The nonlinear solver prints a table per step.  The key columns map directly to
the internal quantities:

\begin{longtable}{p{0.22\linewidth} p{0.68\linewidth}}
\caption{Typical nonlinear output fields.}
\label{tab:output-fields}\\
\toprule
Output field & Meaning \\
\midrule
\endfirsthead
\toprule
Output field & Meaning \\
\midrule
\endhead
\code{time} &
Current target time/load coordinate. \\

\code{timestep} &
Current \(\Delta t\). \\

\code{iter} &
Nonlinear iteration number. \\

\code{residuals} &
\code{resErr}, the normalized force residual error. \\

\code{increments} &
\code{disErr}, the normalized correction/increment error. \\

\code{energies} &
\code{eneErr}, the normalized residual-work error. \\
\bottomrule
\end{longtable}

For iterative linear solvers, the more detailed event file records each linear
system inside each nonlinear iteration.  This is the right place to check
whether more nonlinear work is caused by changed matrix updates, worse
preconditioner quality, loose linear tolerance, or increased linear iterations.

\section{Summary}

The \oas{} nonlinear solver is a load/time-step Newton-like equilibrium solver
with:

\begin{itemize}[leftmargin=2em]
  \item a clear separation between committed state \code{r} and trial state
        \code{trial_r};
  \item full-space force assembly followed by reduced-space linear solves;
  \item explicit constraint and prescribed-boundary handling through
        \code{NodeContainer};
  \item residual, increment, and energy convergence checks separated by
        physical field;
  \item restart and adaptive time-step control when convergence is difficult;
  \item optional relaxed-tolerance acceptance;
  \item optional indirect displacement control using two linear solves per
        nonlinear iteration;
  \item detailed runtime and linear-solver profiling hooks.
\end{itemize}

The central nonlinear equation is the active equilibrium residual
\[
  R_f(r^{\mathrm{tr}}, t) = 0,
\]
solved by repeated linear corrections
\[
  K_{\mathrm{eff}} \Delta r_f = R_f,
\]
where the definitions of \(K_{\mathrm{eff}}\), \(R_f\), and the convergence
scales are determined by the current stiffness update policy, constraints,
physical fields, and selected linear solver backend.

\end{document}
